\clearpage
\SourceLeaf{453}
\begin{center}\large [ZUR THEORIE DES KRÜMMUNGSMAASSES.]\end{center}
\noindent\hspace*{1.2em}Die Bedingungsgleichung, dass
\[
pdx^2+2qdxdy+rdy^2
\]
\noindent das reine Product aus zwei vollständigen Differentialen sei, ist folgende:
\[
\begin{aligned}
0={}&2(qq-pr)\left(\frac{\partial\partial p}{\partial y^2}
-2\frac{\partial\partial q}{\partial x\partial y}
+\frac{\partial\partial r}{\partial x^2}\right)\\
&+p\left(\frac{\partial p}{\partial y}.\frac{\partial r}{\partial y}
-2\frac{\partial q}{\partial x}.\frac{\partial r}{\partial y}
+\frac{\partial r}{\partial x}.\frac{\partial r}{\partial x}\right)\\
&+q\left(\frac{\partial p}{\partial x}.\frac{\partial r}{\partial y}
-\frac{\partial p}{\partial y}.\frac{\partial r}{\partial x}
+4\frac{\partial q}{\partial x}.\frac{\partial q}{\partial y}\right.\\
&\qquad\left.-2\frac{\partial q}{\partial x}.\frac{\partial r}{\partial x}
-2\frac{\partial p}{\partial y}.\frac{\partial q}{\partial y}\right)\\
&+r\left(\frac{\partial p}{\partial x}.\frac{\partial r}{\partial x}
-2\frac{\partial p}{\partial x}.\frac{\partial q}{\partial y}
+\frac{\partial p}{\partial y}.\frac{\partial p}{\partial y}\right).
\end{aligned}
\]

\Needspace{9\baselineskip}
\begin{center}\textbf{BEMERKUNGEN.}\end{center}
\noindent\hspace*{1.2em}Nachdem \textsc{Gauss} den Ausdruck für das Krümmungsmaass einer Fläche zuerst unter der Voraussetzung, dass das Quadrat des Linienelementes die Form
\[
dS^2=mm(dt^2+du^2)
\]
\noindent besitzt (S.\ 381 und 385 dieses Bandes) und darauf unter der Annahme geodätischer Coordinaten, für die
\[
dS^2=ds^2+mmd\theta^2
\]
\noindent wird (S.\ 407 und 442 dieses Bandes), hergeleitet hatte, ist er, wie die vorstehende Notiz zeigt, die sich in einem Handbuche findet, zu dem Zähler des Ausdruckes des Krümmungsmaasses bei \textls[140]{beliebigen} Coordinaten durch die Überlegung gelangt, dass das Krümmungsmass der auf die Ebene abwickelbaren Flächen verschwindet. Die vollständige Formel für das Krümmungsmass wird dann in demselben Handbuche auf dieselbe Art, wie in \S\ 11 der \textit{Disquisitiones generales circa superficies curvas} hergeleitet, sodass es nicht erforderlich schien, diese Rechnungen zu veröffentlichen. Unmittelbar darauf folgt mit derselben Schrift und Tinte die auf Seite 447 bis 448 abgedruckte Notiz: \textit{Allgemeinste Auflösung des Problems der Abwickelung der Flächen.}
\begin{flushright}\textsc{Stäckel}.\end{flushright}

